%==========================================================================
% Addendum to disaster_corrected.tex
%
% Covers three extensions to the canonical model that were added to the
% codebase after disaster_corrected.tex (Feb 2026) was written:
%
%   D1. Effective lower bound (ELB) on the gross policy rate, applied as
%       a high-sharpness softplus floor on the Taylor-rule output.
%   D2. The R_lb_sharpness calibration parameter that controls the kink.
%   D3. Bernoulli disaster mechanism: per-period capital destruction with
%       probability p_disaster, magnitude theta_disaster.
%
% Drop-in: insert these sections after Section 3 (Equilibrium Equations)
% and before Section 4 (Analytical Eliminations) in disaster_corrected.tex.
% Update the calibration table in Section 5 to include the four new
% parameters listed at the end.
%
% Style matches disaster_corrected.tex: color macros for variables, boxed
% canonical formulas, paragraph-level subdivision.
%==========================================================================

%==========================================================================
\section{Effective Lower Bound on the Policy Rate}
\label{sec:elb}
%==========================================================================

The Taylor rule defined in Section~2 (Definitions) prescribes a target
gross nominal rate $R^{\text{Taylor}}_t$ that can fall arbitrarily far
below unity in deep recessions or stress states. Empirically central
banks face a lower bound on the policy rate set by cash-storage costs
and preferences for deposit holdings; historically the bound has been
slightly negative in net terms (the SNB held a policy rate of
${\sim}-0.75\%$ annual for seven years; ECB and BoJ similar) but
typically not deeply so.

We impose this bound on the realized rate $R_t$ via a high-sharpness
softplus floor:
\begin{equation*}
  \boxed{\textcolor{brown}{R_t}
    = R^{lb}
      + \frac{1}{\sharp}\,\log\!\left(
          1 + \exp\!\bigl(\sharp\,(\textcolor{brown}{R^{\text{Taylor}}_t} - R^{lb})\bigr)
        \right)}
\end{equation*}
where $R^{lb}$ is the floor (default $1.0$, a zero net rate) and
$\sharp \equiv R^{lb}_{\text{sharpness}}$ is the kink sharpness
(default $500$).

\paragraph{Properties.}
\begin{itemize}
\item $C^{\infty}$ in $R^{\text{Taylor}}_t$ but with a sharp transition
  near $R^{\text{Taylor}}_t = R^{lb}$. For $R^{\text{Taylor}}_t \gg R^{lb}$,
  $R_t \approx R^{\text{Taylor}}_t$. For $R^{\text{Taylor}}_t \ll R^{lb}$,
  $R_t \approx R^{lb}$.
\item At the steady state, $R^{\text{Taylor}}_{ss} = R_{ss} = 1.018 \gg R^{lb}$,
  so the floor is slack in expectation and the SS distortion from the
  softplus is ${\sim}10^{-7}$.
\item Under stress shocks (especially monetary policy shocks $m_t^p$
  combined with disasters), $R^{\text{Taylor}}_t$ can dip below $R^{lb}$
  and the floor activates.
\end{itemize}

\paragraph{Replaces.}
The unfloored bare-Taylor-rule formula in Section~2:
\begin{equation*}
  \textcolor{brown}{R_t}
    = R_{\text{ss}} \left(\frac{\textcolor{green!60!black}{R_{t-1}}}{R_{\text{ss}}}\right)^{\!\rho_p}
      \!\left[
        \left(\frac{\textcolor{orange!80!black}{\pi_t}}{\pi_{\text{ss}}}\right)^{\!\alpha_\pi}
        \left(\frac{\textcolor{brown}{y_t^{\text{GDP}}}}{y_{\text{ss}}}\right)^{\!\alpha_y}
      \right]^{1-\rho_p}
      \!\exp(\textcolor{green!60!black}{m_t^p})
\end{equation*}
should be renamed to $R^{\text{Taylor}}_t$ (the unfloored prescription)
and the realized $R_t$ defined as above via the softplus floor. The bond
Euler (Eq.~6) and the Taylor-rule definition continue to use $R_t$
unchanged.

\paragraph{Sharpness calibration.}
$\sharp = 500$ is the production default. Lower sharpness gives a
softer transition (more gradient through the kink, but more SS
distortion); higher sharpness sharpens the kink toward the limiting
hard $\max(R^{\text{Taylor}}_t,\, R^{lb})$.

\medskip
\noindent\emph{Implementation note.} The softplus is applied via the
helper $\lfloor x \rfloor_\varepsilon^{\sharp} = \varepsilon +
\frac{1}{\sharp}\log(1 + e^{\sharp(x - \varepsilon)})$ already used for
the Calvo inner terms (Section~2), with $\varepsilon = R^{lb}$ and
$\sharp$ set by configuration.


%==========================================================================
\section{Disaster Mechanism}
\label{sec:disaster}
%==========================================================================

The model admits a Barro-style exogenous disaster shock that destroys a
fraction of next-period capital. Disaster realizations are i.i.d.\ across
time and independent of the continuous Gaussian shocks.

\subsection{Specification}

Each period a Bernoulli disaster indicator is realized:
\begin{equation*}
  d_t \in \{0, 1\}, \qquad
  \mathbb{P}(d_t = 1) = p_{\text{disaster}}
\end{equation*}
with default $p_{\text{disaster}} = 0$ (baseline CMR; mechanism
inactive) and stress calibration $p_{\text{disaster}} = 0.10$
(per-quarter probability).

When $d_t = 1$, the capital entering next period is destroyed by
log-fraction $\theta_{\text{disaster}}$:
\begin{equation*}
  \boxed{\textcolor{brown}{k_t}\bigm|_{d_t}
    = \textcolor{brown}{k_t^{\text{raw}}}
      \cdot \exp(-\theta_{\text{disaster}}\, d_t)}
\end{equation*}
where $k_t^{\text{raw}}$ is the capital implied by the standard law of
motion (Section~2) and $k_t\bigm|_{d_t}$ is the disaster-adjusted
capital that becomes the next-period state.

For the default $\theta_{\text{disaster}} = 0.05$, a disaster destroys
${\sim}4.9\%$ of the capital stock in a single period.

\subsection{Mixture expectation}

All conditional expectations $\mathbb{E}_t[\,\cdot\,]$ in the
forward-looking equilibrium equations (Eq.~1, 2b, 3, 4b, 5, 6, 7, 8 in
Section~3) become mixtures over the disaster indicator:
\begin{equation*}
  \boxed{\mathds{E}_t[X_{t+1}]
    = (1 - p_{\text{disaster}})\,
      \mathds{E}_t\!\left[X_{t+1} \,\big|\, d_{t+1} = 0\right]
    + p_{\text{disaster}}\,
      \mathds{E}_t\!\left[X_{t+1} \,\big|\, d_{t+1} = 1\right]}
\end{equation*}

The two conditional branches differ only through the next-period
state: in the disaster branch ($d_{t+1} = 1$), $k_{t+1}$ is replaced by
$k_{t+1} \cdot e^{-\theta_{\text{disaster}}}$ before the next-period
quantities (definitions, controls, residuals) are computed. The
continuous shocks $\xi_{t+1}$ and exogenous AR(1) transitions are
identical across the two branches.

When $p_{\text{disaster}} = 0$, the mixture reduces to the no-disaster
branch and the model is exactly baseline CMR.

\paragraph{Scope.}
The disaster shock affects only the physical capital stock. It does not
shift the financial-friction parameters ($\sigma_\omega, \mu_{\text{mon}},
\gamma_e, w^e$), the productivity processes, or the monetary/fiscal
parameters. The financial accelerator (Section~2) responds endogenously
to the lower next-period $k_{t+1}$ via $R^k_{t+1}$, $\bar\omega_{t+1}$,
and $L_{t+1}$ in the disaster branch.

\paragraph{Numerical handling.}
Computing the mixture exactly requires evaluating each forward-looking
residual at both disaster realizations and combining with weights
$(1 - p_{\text{disaster}}, p_{\text{disaster}})$. This doubles the
forward-pass cost relative to baseline.


%==========================================================================
\section*{Calibration table additions}
%==========================================================================

The following four parameters should be appended to Table~1 in
Section~5 (Calibration):

\begin{center}
\begin{tabular}{lcll}
\toprule
Parameter & Symbol & Value & Source \\
\midrule
Effective lower bound on gross rate & $R^{lb}$ & $1.0$ & Configuration default \\
ELB softplus sharpness & $R^{lb}_{\text{sharpness}}$ & $500$ & Calibration default \\
Per-period disaster probability & $p_{\text{disaster}}$ & $0.0$--$0.10$ & $0$ baseline; $0.10$ stress \\
Disaster magnitude (log-fraction) & $\theta_{\text{disaster}}$ & $0.05$ & Calibration default \\
\bottomrule
\end{tabular}
\end{center}

\paragraph{Notes.}
\begin{itemize}
\item $R^{lb} = 1.0$ corresponds to a zero net nominal rate. Lowering
  $R^{lb}$ to $\sim 0.998$ would approximate the SNB's $-0.75\%$ annual
  policy rate at quarterly frequency.
\item $R^{lb}_{\text{sharpness}} = 500$ keeps the SS distortion from
  the softplus floor at ${\sim}10^{-7}$ while still producing a tight
  kink near the floor. Lower sharpness ($\sim 100$) gives more
  gradient through the kink; higher sharpness approaches a hard
  $\max(\cdot, R^{lb})$.
\item $p_{\text{disaster}} = 0$ is the production default; the
  mechanism is inert and the model is exactly baseline CMR. Stress
  experiments use $p_{\text{disaster}} = 0.10$.
\item $\theta_{\text{disaster}} = 0.05$ corresponds to a ${\sim}4.9\%$
  capital destruction per disaster realization.
\end{itemize}


%==========================================================================
\section*{Errata addendum}
%==========================================================================

The following extensions should be added to Section~7 (Errata Summary)
of disaster\_corrected.tex:

\begin{enumerate}
\setcounter{enumi}{6}
\item \textbf{ELB floor on $R_t$} (Section~D1 above): the bare Taylor
  rule in the original Section~2 (Definitions) is replaced by a
  softplus-floored realized rate. The unfloored prescription
  $R^{\text{Taylor}}_t$ is computed first; the realized $R_t$ used in
  the bond Euler is the softplus floor of that.
\item \textbf{Disaster mechanism} (Section~D2 above): an i.i.d.\
  Bernoulli disaster shock destroys a fraction $1 - e^{-\theta}$ of
  next-period capital with probability $p_{\text{disaster}}$. All
  forward-looking equilibrium expectations become mixtures over the
  disaster indicator. Default $p_{\text{disaster}} = 0$ recovers
  baseline CMR exactly.
\end{enumerate}
